\section{Validation}
\label{app:validation}

\paragraph{Synthetic preferential attachment.}
We grow synthetic networks of 8,000 nodes in which each new node attaches to an existing node with weight
$d^{\alpha}$ if its degree $d$ is positive and $e^{\beta_0}$ otherwise ($\alpha = 1$, $\beta_0 = -1$), replay them with the same code as the real data and estimate
$(\1[d>0], \log d)$ under each of the three sampling designs. In all three, $\hat\alpha$ lies within three standard
errors of 1. The empirical inclusion frequency of each node under PPS matches $\pi$, and the expected set size
matches $S$.

\paragraph{Recovering simulated mechanisms.}
For \cref{sec:validation} we keep every real candidate set and replace the chosen parent by a draw with probability
proportional to $\exp(V_{kl} - \log \pi_{kl})$, where $V$ contains (a) only the popularity terms of the estimated
model, (b) only quality, demand match and the externally explained part of the hub fixed effects, or (c) the full
model with fixed effects. Re-running the four-group decomposition on these choices gives
\recSfPop\% popularity for (a), \recQfFit\% fit for (b), and \recFullPop\% / \recFullMatch\% / \recFullInertia\%
for (c). Because the candidate features are the real ones, the popularity and match signals keep their real
correlation; (a) and (b) therefore do not reach 100\%.

\paragraph{Full data versus a 10\% sample.}
An earlier version of the analysis used a 10\% random sample of choices. \Cref{tab:validate} compares it with the
full-data estimates under the stratified design (the design of the earlier analysis). \valIn{} of \valN{} coefficients
lie inside the 95\% interval of the 10\% sample. The exceptions are features changed in the rebuild: major
organization (the re-distributor flag now separates accounts that republish other organizations' weights), missing
card text (new text pipeline) and features based on dataset profiles, which are now computed only from earlier models.

\paragraph{Why standard errors shrink by less than $\sqrt{10}$.}
With independent choices the full-data standard errors would be $\sqrt{0.1} \approx 0.32$ of the 10\% ones; the
median ratio is \valSeRatio. Choices are clustered by author: the \nFirst{} first adoptions come from \nFirstAuthors{}
authors, and a 10\% sample of choices keeps 15\% of the authors. Under clustering the variance scales with
$\sum_a n_a^2 / N^2$, which bounds the ratio between 0.32 (independent choices within an author) and 0.89 (perfectly
dependent choices); the observed value implies strong within-author dependence, which is why every interval in the
paper is clustered or bootstrapped by author.

\begin{table}[htbp]
\centering
\caption{Full-data estimates (stratified design) against the earlier 10\% sample. Differences outside the 10\% interval concern features redefined in the rebuild.}
\label{tab:validate}
\scriptsize
\begin{tabular}{lccrr}
\toprule
Feature & 10\% sample & Full & In CI & SE ratio \\
\midrule
\texttt{author\_family} & -0.062 (0.085) & -0.065 (0.062) & yes & 0.73 \\
\texttt{bench\_na} & 0.176 (0.096) & -0.055 (0.048) & no & 0.50 \\
\texttt{bench\_z} & 0.015 (0.080) & -0.110 (0.025) & yes & 0.31 \\
\texttt{ds\_couse} & 1.150 (0.040) & 1.217 (0.020) & yes & 0.48 \\
\texttt{ds\_lang} & 0.459 (0.122) & 0.340 (0.051) & yes & 0.42 \\
\texttt{ds\_task} & 1.245 (0.149) & 0.970 (0.116) & yes & 0.78 \\
\texttt{has\_deg} & 4.687 (0.079) & 4.711 (0.082) & yes & 1.03 \\
\texttt{lang\_hist} & 0.360 (0.042) & 0.239 (0.027) & no & 0.64 \\
\texttt{lco2} & 0.676 (0.135) & 0.773 (0.095) & yes & 0.70 \\
\texttt{legacy} & 1.372 (0.052) & 1.379 (0.033) & yes & 0.64 \\
\texttt{lib\_diffusers} & 0.442 (0.269) & 0.555 (0.154) & yes & 0.57 \\
\texttt{lib\_peft} & -1.763 (0.305) & -1.674 (0.139) & yes & 0.46 \\
\texttt{lib\_sentence-transformers} & 0.873 (0.113) & 0.941 (0.047) & yes & 0.41 \\
\texttt{lib\_transformers} & 0.196 (0.082) & 0.202 (0.036) & yes & 0.45 \\
\texttt{log\_age} & -0.539 (0.021) & -0.540 (0.019) & yes & 0.90 \\
\texttt{log\_deg} & 0.347 (0.014) & 0.336 (0.011) & yes & 0.79 \\
\texttt{lquant} & 0.099 (0.013) & 0.070 (0.008) & no & 0.65 \\
\texttt{lrec30} & 0.631 (0.019) & 0.617 (0.015) & yes & 0.80 \\
\texttt{major\_org} & 0.300 (0.046) & 0.681 (0.031) & no & 0.68 \\
\texttt{same\_author} & 6.566 (0.449) & 6.319 (0.504) & yes & 1.12 \\
\texttt{size\_dist} & -0.173 (0.025) & -0.226 (0.017) & no & 0.66 \\
\texttt{size\_na} & -0.780 (0.062) & -0.824 (0.035) & yes & 0.56 \\
\texttt{task\_hist} & 1.909 (0.077) & 1.882 (0.044) & yes & 0.57 \\
\texttt{txt\_na} & -0.743 (0.074) & -0.368 (0.058) & no & 0.78 \\
\texttt{txt\_sim} & 4.559 (0.656) & 4.360 (0.757) & yes & 1.15 \\
\bottomrule
\end{tabular}
\end{table}

